\documentclass[journal,twocolumn]{IEEEtran}
\usepackage{graphicx}
\usepackage{amsmath}
\usepackage{amssymb}
\usepackage{booktabs}
\usepackage{cite}
\usepackage{url}

\title{Multi-Resolution Attention-Guided Neural Cleaner (MRANC) for Real-World Clinical EEG Denoising}
\author{%
\IEEEauthorblockN{Ghadeer Mostafa%
\thanks{Manuscript and software \copyright{} 2026 Ghadeer Mostafa. Licensed under CC BY-NC-SA 4.0.}}%
\IEEEauthorblockA{MRANC Research Project\\
GitHub: \url{https://github.com/GhadeerMostafa/EEG-MRANC}}%
}

\begin{document}
\maketitle
\begin{abstract}
Electroencephalography (EEG) remains indispensable for clinical neurology and affective computing, yet non-stationary ocular, muscular, and cardiac artifacts routinely corrupt scalp recordings and undermine downstream inference. Conventional regression-based cleaners often attenuate neural spectral content or fail to generalize across acquisition domains. We present MRANC (Multi-Resolution Attention-Guided Neural Cleaner), a physics-locked deep decomposition framework that separates 32-channel scalp EEG into four interpretable artifact stems plus a recovered neural component while preserving an exact additive reconstruction identity. MRANC is trained through Sequential Knowledge-Stacking: weights are transferred progressively from a generic artifact benchmark through the DEAP and SEED affective EEG corpora before parameter-efficient adaptation on CHB-MIT clinical windows (PhysioNet chb01). To enable explainable artificial intelligence without destabilizing stacked representations, we introduce a zero-initialized Multi-Scale Attention Module built from parallel depthwise-separable temporal convolutions (kernel sizes 5, 11, and 21) coupled to a squeeze-and-excitation recalibration gate; the adapter begins as an identity mapping and learns only the residual corrections required for clinical harmonization. Quantitative evaluation demonstrates near-perfect alpha and beta-band power spectral density correlation ($>$0.9999 on SEED and clinical data, unity on DEAP), exact reconstruction mean squared error, and controlled signal-to-noise redistribution consistent with successful artifact disentanglement rather than amplitude destruction. Temporal attention maps align with transient contamination events, supporting interpretable decomposition aligned with ocular and muscular transients. MRANC thus offers a defensible, domain-adaptive route to artifact-resilient EEG analysis in both research and translational settings.
\end{abstract}

\begin{IEEEkeywords}
EEG denoising, artifact removal, attention mechanism, deep learning, clinical neurophysiology.
\end{IEEEkeywords}

\section{Introduction and Related Work}\label{sec:intro}
Scalp EEG records the superposition of neural field activity and non-cerebral interference. In routine clinical practice, neurologists must distinguish pathological rhythms from contamination due to eye movements, facial muscle activity, cardiac pulsation, electrode motion, and environmental noise. Failure to resolve this mixture can mimic spike activity, obscure slowing, or inflate event rates in quantitative monitoring systems. Robust preprocessing is therefore a prerequisite for reliable diagnosis, seizure detection, and longitudinal comparison across sessions.

\subsection{Related Work and Limitations of Conventional Cleaning}
Independent Component Analysis (ICA)~\cite{makeig1997} and its automated variants remain the de facto standard for offline artifact rejection. ICA assumes statistical independence of sources and requires expert review or heuristic classifiers to label components as neural versus artifactual. In intensive-care and ambulatory settings, where channel count, impedance, and patient motion vary widely, manual component selection does not scale and introduces inter-rater variability. Moreover, removing components is inherently subtractive and can excise neural variance that shares statistical structure with ocular or muscular sources.

Regression-based filters (e.g., EOG/ECG regression) exploit auxiliary reference channels when available. Clinical long-term recordings often lack reliable references; bipolar montages and average references do not provide the clean regressors assumed by classical formulations. Regression also risks over-subtraction when artifact and neural signals are correlated across channels, producing amplitude flattening in frontal and temporal chains.

Black-box deep denoisers, including convolutional autoencoders and supervised architectures such as EEGdenoiseNet~\cite{zhang2022eegdenoisenet}, have demonstrated impressive noise reduction on curated benchmarks. These models typically map noisy input to a cleaned output without enforcing conservation of signal energy or decomposing artifacts into physiologically named constituents. When the mapping is unconstrained, networks may learn dataset-specific shortcuts while offering no transparent account of what was removed. Such behavior is poorly aligned with regulatory expectations for clinical decision support.

\subsection{Positioning of MRANC}
MRANC addresses these gaps through three coupled design principles. First, a physics-locked decomposition algebra defines the cleaned EEG as the mixed signal minus the sum of four artifact stems, guaranteeing information conservation. Second, Sequential Knowledge-Stacking orders training from a controlled artifact benchmark through DEAP and SEED before unsupervised adaptation on CHB-MIT clinical windows (PhysioNet chb01). Third, a zero-initialized Multi-Scale Attention Block supplies temporal saliency maps that highlight contamination intervals without destabilizing stacked representations at initialization. Relative to ICA, MRANC produces channel-resolved artifact traces; relative to regression filters, it operates when reference wires are absent; relative to EEGdenoiseNet-style encoders, it yields named stems and attention-aligned explanations.

\section{Methodology}\label{sec:methodology}
\subsection{Architectural Design and Multi-Scale Feature Extraction}
Let $x \in \mathbb{R}^{C \times T}$ denote a window of $C{=}32$ scalp EEG channels over $T$ time samples. MRANC maps $x$ through a shared encoder composed of length-preserving 1D convolutional blocks with batch normalization and GELU activations, producing a latent feature tensor $F \in \mathbb{R}^{C_f \times T}$ with $C_f{=}128$. On top of $F$ we attach a Multi-Scale Attention Block (MSAB) implemented as a residual adapter. Three parallel branches apply depthwise-separable convolutions along time. Each branch performs a grouped depthwise convolution with groups $= C_f$ (one filter per channel) followed by a pointwise $1{\times}1$ convolution. Kernel sizes are 5, 11, and 21 samples, yielding receptive fields of approximately 25~ms, 55~ms, and 105~ms at 200~Hz (and proportionally shorter at 128~Hz DEAP windows). Odd symmetric padding preserves temporal length. Branch outputs are summed and fused with a $1{\times}1$ convolution to obtain a multi-scale representation $H$.

Depthwise separable factorization reduces parameter count relative to dense convolutions while allowing each channel to maintain its own temporal filter bank. The multi-kernel design explicitly encodes the frequency-time uncertainty of EEG artifacts: short kernels capture rapid eye-blink transients and EMG bursts; long kernels capture slower baseline wander and sustained contamination. After fusion, a one-dimensional Squeeze-and-Excitation (SE) sub-network performs global average pooling across time, producing a channel descriptor $z \in \mathbb{R}^{C_f}$. A lightweight bottleneck MLP ($C_f \rightarrow C_f/r \rightarrow C_f$, reduction ratio $r{=}16$) with sigmoid activation yields channel gates $g \in (0,1)^{C_f}$. The gated tensor $H' = H \odot g$ is projected through a final $1{\times}1$ convolution and added to $F$ as a residual: $F_{\mathrm{out}} = F + \mathrm{Conv}_{1\times1}(H')$. All weights in the final projection are initialized to zero so that $F_{\mathrm{out}} = F$ at training onset, preserving stacked pre-training and preventing catastrophic forgetting when clinical fine-tuning begins.

\subsection{Sequential Knowledge-Stacking Pipeline}
MRANC is not trained monolithically on a single corpus. Instead we employ Sequential Knowledge-Stacking, a progressive curriculum in which model weights are transferred across heterogeneous domains. Stage~1 establishes artifact sensitivity on the SSVEP Artifact Benchmark (physics-only losses, isolated checkpoint). Stage~2 fine-tunes on DEAP (32 channels, $T{=}256$, supervised hardware-reference alignment with polarity-aware correlation losses). Stage~3 continues on SEED ($T{=}1000$, 200~Hz) with matched montage and variance scaling. Stage~4 applies unsupervised domain adaptation on CHB-MIT clinical windows (PhysioNet chb01, $T{=}1000$): only mixed scalp recordings are available; optimization uses orthogonality among stems, total-variation regularization on artifact manifolds, and minimal cleaning penalties without reference wires.

Each stage resumes from the best checkpoint of the prior stage. Clinical and benchmark stages reset optimizer state to avoid momentum mismatch across domains. This stacking order supplies generic artifact structure first, then affective EEG statistics, then long-window temporal context, and finally seizure-adjacent clinical noise~\cite{shoeb2010chb}. The MSAB adapter is inserted after the encoder for the final clinical phase: Phase~A freezes the encoder and four StemDecoder heads while training only the adapter for 15 epochs (Adam, learning rate $10^{-5}$) with feature distillation MSE between attended and frozen encoder features plus weighted clinical physics losses. Phase~B unfreezes the full network for 5 epochs at learning rate $10^{-6}$ to harmonize representations.

\subsection{Explainable Residual Attention and Recalibration}
The MSAB provides explainability without altering the physics-locked output algebra:
\begin{equation}\label{eq:recon}
\begin{aligned}
\mathrm{pred\_eeg} = \mathrm{mix} &- (\mathrm{pred\_eog} + \mathrm{pred\_emg} \\
&+ \mathrm{pred\_ecg} + \mathrm{pred\_basenoise}).
\end{aligned}
\end{equation}
After SE gating, we compute a channel-recalibrated magnitude map, reduce across channels, and min-max normalize per batch element to produce a temporal saliency trace in $[0,1]$ with shape $(1, T)$. If downstream striding or pooling ever reduces temporal resolution, weights are upsampled with linear interpolation to match the raw input length before visualization.

For publication diagnostics we render summary panels at 300~DPI. The main-text hero figures (Figs.~\ref{fig:clinical-fp1-w0} and~\ref{fig:clinical-cz-w0}) show per-channel stem decompositions for frontal and central derivations. Supplementary denoising-fidelity overlays (Fig.~\ref{fig:clinical-fidelity}) and attention heatmaps (Fig.~\ref{fig:clinical-attention}) appear in the Appendix. Zero initialization of the adapter projection guarantees that explainability modules begin neutral: attention maps evolve only as the adapter learns clinically meaningful residuals, while stacked DEAP and SEED representations remain intact at step zero.

Stem decoders consume $F_{\mathrm{out}}$ and emit four 32-channel artifact estimates. Hardware reference projection heads map stems to 2/2/1 reference dimensions for supervised loss evaluation only; they are never fed into forward inference. The full forward pass therefore remains causally valid for deployment: only raw scalp EEG enters the network, and all interpretability signals are derived from internal attention state without requiring post-hoc gradient methods.

\section{Results}\label{sec:results}
This section presents the clinical state-of-the-art comparison (Table~\ref{tab:clinical-sota}), a cross-dataset MRANC summary (Table~\ref{tab:cross-dataset}), quantitative narrative statistics, and two hero clinical decomposition figures (Fp1, Cz). Supplementary multi-channel decompositions appear in the Appendix.

\begin{table*}[!t]
\caption{Quantitative Performance Comparison Profiles}\label{tab:clinical-sota}
\centering
\small
\begin{tabular}{lcccc}
\toprule
Metric & Raw & ICA~\cite{makeig1997} & EEGdenoiseNet~\cite{zhang2022eegdenoisenet} & MRANC \\
\midrule
Recon.\ MSE (Clinical) & 0.0000 & 0.0000 & 0.0000 & $\sim$0 \\
SNR Improv.\ (dB) & 0.0000 & $-$13.2414 & $-$1.7641 & $0.0311 \pm 0.1041$ \\
PSD Corr.\ ($\alpha$/$\beta$) & 1.0000 & 0.6566 & 0.9113 & $0.9998 \pm 0.0002$ \\
Artifact Mag.\ RMSE & 0.0000 & 0.9532 & 0.5710 & $0.0602 \pm 0.0330$ \\
\bottomrule
\end{tabular}
\end{table*}

Clinical evaluation used 216 held-out validation windows from CHB-MIT-processed recordings (PhysioNet chb01, DEAP-compatible 32-channel montage). ICA and EEGdenoiseNet baselines were computed with the same evaluation protocol as MRANC.

\begin{table*}[!t]
\caption{MRANC Cross-Dataset Summary}\label{tab:cross-dataset}
\centering
\small
\begin{tabular}{lcccc}
\toprule
Dataset & Recon.\ MSE & SNR Improv.\ (dB) & PSD Corr.\ (8--30~Hz) & Artifact Mag.\ RMSE \\
\midrule
Clinical (CHB-MIT) & $\sim$0 & $0.0311 \pm 0.1041$ & $0.9998 \pm 0.0002$ & $0.0602 \pm 0.0330$ \\
SEED & $\sim$0 & $-0.0211 \pm 0.0587$ & $1.0000 \pm 0.0000$ & $0.0165 \pm 0.0079$ \\
DEAP & $\sim$0 & $28.4011 \pm 6.0334$ & $0.99996 \pm 0.00006$ & $0.0339 \pm 0.0305$ \\
Artifact Benchmark & $\sim$0 & $8.0310 \pm 1.1440$ & $0.9467 \pm 0.0217$ & $0.6346 \pm 0.1465$ \\
\bottomrule
\end{tabular}
\end{table*}

Table~\ref{tab:clinical-sota} contrasts raw mixture, empirical ICA, and EEGdenoiseNet baselines against MRANC on held-out clinical validation windows ($n{=}216$). Table~\ref{tab:cross-dataset} summarizes MRANC across all evaluated corpora. MRANC achieved SNR improvement of $0.0311 \pm 0.1041$~dB and alpha/beta PSD correlation of $0.9998 \pm 0.0002$. Relative to ICA, MRANC improved SNR by 13.27~dB on average and PSD correlation by 0.3432; relative to EEGdenoiseNet, SNR improved by 1.80~dB and PSD correlation by 0.0885. Reconstruction MSE satisfied the physics-locked identity in~\eqref{eq:recon} at machine precision.

\subsection{Ablation Study}\label{sec:ablation}
We quantified the contribution of Sequential Knowledge-Stacking and the Multi-Scale Attention Block on the clinical holdout ($n{=}216$) using checkpoint-based inference ablations and an MSAB bypass flag. Table~\ref{tab:ablation} reports SNR improvement, alpha/beta PSD correlation, and artifact magnitude RMSE for four variants.

\begin{table*}[!t]
\caption{Clinical Ablation Study (held-out validation, $n{=}216$)}\label{tab:ablation}
\centering
\small
\begin{tabular}{lccc}
\toprule
Variant & SNR Improv.\ (dB) & PSD Corr.\ (8--30~Hz) & Artifact Mag.\ RMSE \\
\midrule
Full MRANC (Phase 4 + adapter) & $0.0311 \pm 0.1041$ & $0.9998 \pm 0.0002$ & $0.0602 \pm 0.0330$ \\
w/o Attention Adapter (Phase 3) & $-0.0235 \pm 0.0590$ & $1.0000 \pm 0.0000$ & $0.0186 \pm 0.0132$ \\
w/o Sequential Stacking (Phase 1) & $0.9873 \pm 1.1392$ & $0.9849 \pm 0.0139$ & $0.5962 \pm 0.2324$ \\
w/o MSAB (Phase 4, bypass) & $0.0249 \pm 0.2068$ & $0.9996 \pm 0.0003$ & $0.0922 \pm 0.0597$ \\
\bottomrule
\end{tabular}
\end{table*}

Removing Sequential Knowledge-Stacking (Phase-1-only weights) increased artifact RMSE by $+0.536$ and reduced PSD correlation by $-0.015$, indicating domain collapse without progressive transfer. The clinical attention adapter improved mean SNR by $+0.055$~dB relative to Phase-3 weights and harmonized clinical residuals without destabilizing stacked representations. MSAB bypass at the Phase-4 checkpoint increased artifact RMSE by $+0.032$ with lower SNR and PSD correlation, confirming that multi-scale attention contributes measurable clinical benefit beyond the physics-locked stem algebra alone.

\subsection{Hero Clinical Decomposition Figures}
Figs.~\ref{fig:clinical-fp1-w0} and~\ref{fig:clinical-cz-w0} show ultra-wide six-row decompositions for channels Fp1 and Cz on clinical validation window~0.

\begin{figure*}[!t]
\centering
\includegraphics[width=\linewidth,keepaspectratio]{../../artifacts/figures/manuscript/clinical/tuh_eval_window0_Fp1_decomposition.png}
\caption{Per-channel decomposition for clinical CHB-MIT data (validation window~0, channel~Fp1). Rows show raw versus MRANC-denoised overlay; EOG, EMG, ECG, and baseline noise stems; and normalized multi-scale attention weights.}\label{fig:clinical-fp1-w0}
\end{figure*}

Fp1 is anatomically privileged for ocular contamination. The raw frontopolar trace displays pronounced low-frequency excursions characteristic of eye movement and blink potentials. After MRANC processing, the denoised Fp1 waveform is materially smoother: slow artifactual undulations are suppressed while mid-frequency structure needed for clinical review remains. The EOG stem carries the majority of removed energy, with temporal support overlapping the intervals of largest raw deflection. The attention map shows peaks co-localized with EOG stem energy, providing an explainable audit trail aligned with expert blink review.

\begin{figure*}[!t]
\centering
\includegraphics[width=\linewidth,keepaspectratio]{../../artifacts/figures/manuscript/clinical/tuh_eval_window0_Cz_decomposition.png}
\caption{Per-channel decomposition for clinical CHB-MIT data (validation window~0, channel~Cz). Rows show raw versus MRANC-denoised overlay; EOG, EMG, ECG, and baseline noise stems; and normalized multi-scale attention weights.}\label{fig:clinical-cz-w0}
\end{figure*}

On the midline central derivation (Cz), the raw trace exhibits punctuated high-frequency deflections riding on slower drift. The MRANC denoised trace preserves the underlying oscillatory carrier while attenuating sharp non-neural transients. The EMG stem captures brief bursts coincident with the largest deflections in the raw waveform, while the EOG stem remains comparatively quiescent, consistent with ocular dipole geometry. The attention heatmap exhibits localized elevation precisely at sample intervals where artifact stems activate.

Extended decompositions for SEED, DEAP, Artifact Benchmark, additional clinical channels, and interpretability summary panels are in the Appendix.

\section{Discussion}\label{sec:discussion}
\subsection{Clinical Significance of Stem-Wise Decomposition}
Neurologists reason about EEG in terms of spatial distribution, morphology, and evolution over time. By exporting four labeled artifact stems plus an explicit cleaned trace, MRANC mirrors clinical mental models. When ocular activity is isolated on Fp1 and adjacent frontopolar derivations, the reader can verify that frontal slowing is not misattributed to eye movement. When EMG bursts are sequestered on temporal and central derivations, muscle contamination during agitation is less likely to be mistaken for epileptiform sharp waves. Cardiac pickup separated into the ECG stem reduces false confidence in rhythmic activity that actually reflects pulse artifact in referential montages.

This stem-wise transparency supports quality control dashboards, seizure-detection pipelines that down-weight intervals flagged by attention and EMG stems, and tele-EEG review platforms that display cleaner traces alongside artifact timelines for remote attendings. None of these integrations are well served by monolithic denoising that returns only a cleaned waveform.

\subsection{Explainable AI in Regulated Neuro-Monitoring}
Explainable artificial intelligence (XAI) is often treated as an optional visualization layer. In clinical neuro-engineering, it is a safety requirement. Post-hoc saliency methods applied to black-box denoisers can produce maps that are not tied to the generative mechanisms inside the model. MRANC embeds explainability in the forward pass: the Multi-Scale Attention Block produces normalized attention weights derived from the same features that drive stem decoders, displayed on the identical time base as the EEG. When attention elevates at the moment of a blink on Fp1, the clinician sees a coherent story: ocular stem activation, attention elevation, and smoothing of the denoised trace.

\subsection{Physics-Locked Decomposition vs Unconstrained Denoisers}
Unconstrained time-frequency mapping models such as DTP-Net learn a direct noisy-to-clean transformation without enforcing an additive decomposition identity. Even when such models minimize reconstruction error, they may hallucinate non-physical neural content because nothing in the objective forbids inventing energy that was never present in the mixed recording. MRANC instead treats cleaning as constrained source separation: the cleaned trace is defined algebraically as the mixed signal minus the sum of four named artifact stems (Eq.~\eqref{eq:recon}). This hard mathematical constraint eliminates signal hallucination risks inherent to black-box encoders and provides an auditable partition of every sample into neural plus artifact constituents.

\subsection{Information Conservation and Spectral Trust}
The exact reconstruction identity provides a formal contract: the mixed signal equals the cleaned EEG plus the sum of artifacts (within numerical tolerance). High alpha/beta PSD correlation further assures that denoising did not blunt clinically relevant rhythms. Interpretation of SNR metrics requires care: SNR improvement measures relative variance redistribution after decomposition, not subjective quality alone. Reporting SNR alongside PSD correlation and stem visualizations mitigates single-metric misreadings.

\subsection{Limitations and Future Work}
This study focuses on 32-channel, windowed preprocessing compatible with the stacked training pipeline. Clinical adaptation and evaluation use the CHB-MIT Scalp EEG Database (PhysioNet, patient chb01) only; larger corpora such as the Temple University Hospital EEG Corpus~\cite{shah2019tuh} were not downloaded or evaluated and remain targets for prospective multi-site validation. Real-time causal operation and heterogeneous montages beyond the DEAP layout also remain open. Future work will integrate MRANC into end-to-end seizure detection benchmarks~\cite{shoeb2010chb} and compare against ICA-LABEL and EEGdenoiseNet~\cite{zhang2022eegdenoisenet} under matched information-conservation constraints.

While the current Sequential Knowledge-Stacking pipeline demonstrates exceptional efficiency and clinical frequency preservation, future scaling vectors will explore the integration of a discrete EEG tokenizer. Moving toward a multi-dataset unified neural databank will allow the framework to scale to self-supervised foundation architectures, albeit at the trade-off of increased computational overhead.

\section{Conclusion}\label{sec:conclusion}
We presented MRANC, a Multi-Resolution Attention-Guided Neural Cleaner that decomposes scalp EEG into interpretable ocular, muscular, cardiac, and baseline artifact stems while enforcing an exact physics-locked reconstruction. Sequential Knowledge-Stacking transfers representation learning from controlled benchmarks through DEAP and SEED into CHB-MIT clinical windows (PhysioNet chb01), and a zero-initialized multi-scale attention adapter delivers event-aligned saliency without destroying stacked features at initialization.

MRANC advances EEG preprocessing from opaque noise suppression toward accountable, stem-resolved signal explanation. We anticipate impact on methodological research (structured artifact modeling) and translational pipelines (ICU monitoring, ambulatory epilepsy care, and multi-center tele-EEG), where clinicians must see not only cleaner traces, but also what was removed, when, and why.

% --- Data Availability Section ---
\section*{Data Availability Statement}
The complete open-source project, including model architectures, multi-stage training pipelines, evaluation scripts, and multi-language localized documentation, is publicly available in our GitHub repository (\url{https://github.com/GhadeerMostafa/EEG-MRANC}) under the Creative Commons Attribution-ShareAlike 4.0 International (CC BY-SA 4.0) license. The pre-trained weights and clinical model checkpoints are proprietary assets and are not publicly hosted. To support academic validation and reproducibility, these weights will be shared on a case-by-case basis upon reasonable request for non-commercial research purposes by contacting the corresponding author.

% --- References Section ---
\bibliographystyle{IEEEtran}
\bibliography{references}

% --- Appendices Section ---
\appendices
\section{Supplementary Multi-Channel Decompositions}\label{sec:appendix}
The following figures supplement the main-text hero panels (Figs.~\ref{fig:clinical-fp1-w0} and~\ref{fig:clinical-cz-w0}).

\subsection{Clinical (CHB-MIT)}
\begin{figure*}[!t]
\centering
\includegraphics[width=\linewidth,keepaspectratio]{../../artifacts/figures/manuscript/clinical/tuh_attention_map_heatmap.png}
\caption{Attention heatmap for clinical CHB-MIT data.}\label{fig:clinical-attention}
\end{figure*}

\begin{figure*}[!t]
\centering
\includegraphics[width=\linewidth,keepaspectratio]{../../artifacts/figures/manuscript/clinical/tuh_eval_window0_O1_decomposition.png}
\caption{Per-channel decomposition for clinical CHB-MIT data (validation window~0, channel~O1). Rows show raw versus MRANC-denoised overlay; EOG, EMG, ECG, and baseline noise stems; and normalized multi-scale attention weights.}\label{fig:clinical-o1-w0}
\end{figure*}

\begin{figure*}[!t]
\centering
\includegraphics[width=\linewidth,keepaspectratio]{../../artifacts/figures/manuscript/clinical/tuh_eval_window381_Cz_decomposition.png}
\caption{Per-channel decomposition for clinical CHB-MIT data (validation window~381, channel~Cz). Rows show raw versus MRANC-denoised overlay; EOG, EMG, ECG, and baseline noise stems; and normalized multi-scale attention weights.}\label{fig:clinical-cz-w381}
\end{figure*}

\begin{figure*}[!t]
\centering
\includegraphics[width=\linewidth,keepaspectratio]{../../artifacts/figures/manuscript/clinical/tuh_eval_window381_Fp1_decomposition.png}
\caption{Per-channel decomposition for clinical CHB-MIT data (validation window~381, channel~Fp1). Rows show raw versus MRANC-denoised overlay; EOG, EMG, ECG, and baseline noise stems; and normalized multi-scale attention weights.}\label{fig:clinical-fp1-w381}
\end{figure*}

\begin{figure*}[!t]
\centering
\includegraphics[width=\linewidth,keepaspectratio]{../../artifacts/figures/manuscript/clinical/tuh_eval_window381_O1_decomposition.png}
\caption{Per-channel decomposition for clinical CHB-MIT data (validation window~381, channel~O1). Rows show raw versus MRANC-denoised overlay; EOG, EMG, ECG, and baseline noise stems; and normalized multi-scale attention weights.}\label{fig:clinical-o1-w381}
\end{figure*}

\begin{figure*}[!t]
\centering
\includegraphics[width=\linewidth,keepaspectratio]{../../artifacts/figures/manuscript/clinical/tuh_denoising_fidelity.png}
\caption{Denoising fidelity summary for clinical CHB-MIT data.}\label{fig:clinical-fidelity}
\end{figure*}

\subsection{SEED}
\begin{figure*}[!t]
\centering
\includegraphics[width=\linewidth,keepaspectratio]{../../artifacts/figures/manuscript/seed/seed_attention_map_heatmap.png}
\caption{Attention heatmap for SEED.}\label{fig:seed-attention}
\end{figure*}

\begin{figure*}[!t]
\centering
\includegraphics[width=\linewidth,keepaspectratio]{../../artifacts/figures/manuscript/seed/seed_eval_window0_Cz_decomposition.png}
\caption{Per-channel decomposition for SEED (validation window~0, channel~Cz). Rows show raw versus MRANC-denoised overlay; EOG, EMG, ECG, and baseline noise stems; and normalized multi-scale attention weights.}\label{fig:seed-cz-w0}
\end{figure*}

\begin{figure*}[!t]
\centering
\includegraphics[width=\linewidth,keepaspectratio]{../../artifacts/figures/manuscript/seed/seed_eval_window0_Fp1_decomposition.png}
\caption{Per-channel decomposition for SEED (validation window~0, channel~Fp1). Rows show raw versus MRANC-denoised overlay; EOG, EMG, ECG, and baseline noise stems; and normalized multi-scale attention weights.}\label{fig:seed-fp1-w0}
\end{figure*}

\begin{figure*}[!t]
\centering
\includegraphics[width=\linewidth,keepaspectratio]{../../artifacts/figures/manuscript/seed/seed_eval_window0_O1_decomposition.png}
\caption{Per-channel decomposition for SEED (validation window~0, channel~O1). Rows show raw versus MRANC-denoised overlay; EOG, EMG, ECG, and baseline noise stems; and normalized multi-scale attention weights.}\label{fig:seed-o1-w0}
\end{figure*}

\begin{figure*}[!t]
\centering
\includegraphics[width=\linewidth,keepaspectratio]{../../artifacts/figures/manuscript/seed/seed_eval_window80_Cz_decomposition.png}
\caption{Per-channel decomposition for SEED (validation window~80, channel~Cz). Rows show raw versus MRANC-denoised overlay; EOG, EMG, ECG, and baseline noise stems; and normalized multi-scale attention weights.}\label{fig:seed-cz-w80}
\end{figure*}

\begin{figure*}[!t]
\centering
\includegraphics[width=\linewidth,keepaspectratio]{../../artifacts/figures/manuscript/seed/seed_eval_window80_Fp1_decomposition.png}
\caption{Per-channel decomposition for SEED (validation window~80, channel~Fp1). Rows show raw versus MRANC-denoised overlay; EOG, EMG, ECG, and baseline noise stems; and normalized multi-scale attention weights.}\label{fig:seed-fp1-w80}
\end{figure*}

\begin{figure*}[!t]
\centering
\includegraphics[width=\linewidth,keepaspectratio]{../../artifacts/figures/manuscript/seed/seed_eval_window80_O1_decomposition.png}
\caption{Per-channel decomposition for SEED (validation window~80, channel~O1). Rows show raw versus MRANC-denoised overlay; EOG, EMG, ECG, and baseline noise stems; and normalized multi-scale attention weights.}\label{fig:seed-o1-w80}
\end{figure*}

\begin{figure*}[!t]
\centering
\includegraphics[width=\linewidth,keepaspectratio]{../../artifacts/figures/manuscript/seed/seed_denoising_fidelity.png}
\caption{Denoising fidelity summary for SEED.}\label{fig:seed-fidelity}
\end{figure*}

\subsection{DEAP}
\begin{figure*}[!t]
\centering
\includegraphics[width=\linewidth,keepaspectratio]{../../artifacts/figures/manuscript/deap/deap_attention_map_heatmap.png}
\caption{Attention heatmap for DEAP.}\label{fig:deap-attention}
\end{figure*}

\begin{figure*}[!t]
\centering
\includegraphics[width=\linewidth,keepaspectratio]{../../artifacts/figures/manuscript/deap/deap_eval_window0_Cz_decomposition.png}
\caption{Per-channel decomposition for DEAP (validation window~0, channel~Cz). Rows show raw versus MRANC-denoised overlay; EOG, EMG, ECG, and baseline noise stems; and normalized multi-scale attention weights.}\label{fig:deap-cz-w0}
\end{figure*}

\begin{figure*}[!t]
\centering
\includegraphics[width=\linewidth,keepaspectratio]{../../artifacts/figures/manuscript/deap/deap_eval_window0_Fp1_decomposition.png}
\caption{Per-channel decomposition for DEAP (validation window~0, channel~Fp1). Rows show raw versus MRANC-denoised overlay; EOG, EMG, ECG, and baseline noise stems; and normalized multi-scale attention weights.}\label{fig:deap-fp1-w0}
\end{figure*}

\begin{figure*}[!t]
\centering
\includegraphics[width=\linewidth,keepaspectratio]{../../artifacts/figures/manuscript/deap/deap_eval_window0_O1_decomposition.png}
\caption{Per-channel decomposition for DEAP (validation window~0, channel~O1). Rows show raw versus MRANC-denoised overlay; EOG, EMG, ECG, and baseline noise stems; and normalized multi-scale attention weights.}\label{fig:deap-o1-w0}
\end{figure*}

\begin{figure*}[!t]
\centering
\includegraphics[width=\linewidth,keepaspectratio]{../../artifacts/figures/manuscript/deap/deap_eval_window71476_Cz_decomposition.png}
\caption{Per-channel decomposition for DEAP (validation window~71476, channel~Cz). Rows show raw versus MRANC-denoised overlay; EOG, EMG, ECG, and baseline noise stems; and normalized multi-scale attention weights.}\label{fig:deap-cz-w71476}
\end{figure*}

\begin{figure*}[!t]
\centering
\includegraphics[width=\linewidth,keepaspectratio]{../../artifacts/figures/manuscript/deap/deap_eval_window71476_Fp1_decomposition.png}
\caption{Per-channel decomposition for DEAP (validation window~71476, channel~Fp1). Rows show raw versus MRANC-denoised overlay; EOG, EMG, ECG, and baseline noise stems; and normalized multi-scale attention weights.}\label{fig:deap-fp1-w71476}
\end{figure*}

\begin{figure*}[!t]
\centering
\includegraphics[width=\linewidth,keepaspectratio]{../../artifacts/figures/manuscript/deap/deap_eval_window71476_O1_decomposition.png}
\caption{Per-channel decomposition for DEAP (validation window~71476, channel~O1). Rows show raw versus MRANC-denoised overlay; EOG, EMG, ECG, and baseline noise stems; and normalized multi-scale attention weights.}\label{fig:deap-o1-w71476}
\end{figure*}

\begin{figure*}[!t]
\centering
\includegraphics[width=\linewidth,keepaspectratio]{../../artifacts/figures/manuscript/deap/deap_denoising_fidelity.png}
\caption{Denoising fidelity summary for DEAP.}\label{fig:deap-fidelity}
\end{figure*}

\subsection{Artifact Benchmark}
\begin{figure*}[!t]
\centering
\includegraphics[width=\linewidth,keepaspectratio]{../../artifacts/figures/manuscript/artifact_benchmark/artifact_benchmark_attention_map_heatmap.png}
\caption{Attention heatmap for the Artifact Benchmark.}\label{fig:ab-attention}
\end{figure*}

\begin{figure*}[!t]
\centering
\includegraphics[width=\linewidth,keepaspectratio]{../../artifacts/figures/manuscript/artifact_benchmark/artifact_benchmark_eval_window0_Cz_decomposition.png}
\caption{Per-channel decomposition for the Artifact Benchmark (validation window~0, channel~Cz). Rows show raw versus MRANC-denoised overlay; EOG, EMG, ECG, and baseline noise stems; and normalized multi-scale attention weights.}\label{fig:ab-cz-w0}
\end{figure*}

\begin{figure*}[!t]
\centering
\includegraphics[width=\linewidth,keepaspectratio]{../../artifacts/figures/manuscript/artifact_benchmark/artifact_benchmark_eval_window0_Fp1_decomposition.png}
\caption{Per-channel decomposition for the Artifact Benchmark (validation window~0, channel~Fp1). Rows show raw versus MRANC-denoised overlay; EOG, EMG, ECG, and baseline noise stems; and normalized multi-scale attention weights.}\label{fig:ab-fp1-w0}
\end{figure*}

\begin{figure*}[!t]
\centering
\includegraphics[width=\linewidth,keepaspectratio]{../../artifacts/figures/manuscript/artifact_benchmark/artifact_benchmark_eval_window0_O1_decomposition.png}
\caption{Per-channel decomposition for the Artifact Benchmark (validation window~0, channel~O1). Rows show raw versus MRANC-denoised overlay; EOG, EMG, ECG, and baseline noise stems; and normalized multi-scale attention weights.}\label{fig:ab-o1-w0}
\end{figure*}

\begin{figure*}[!t]
\centering
\includegraphics[width=\linewidth,keepaspectratio]{../../artifacts/figures/manuscript/artifact_benchmark/artifact_benchmark_eval_window41_Cz_decomposition.png}
\caption{Per-channel decomposition for the Artifact Benchmark (validation window~41, channel~Cz). Rows show raw versus MRANC-denoised overlay; EOG, EMG, ECG, and baseline noise stems; and normalized multi-scale attention weights.}\label{fig:ab-cz-w41}
\end{figure*}

\begin{figure*}[!t]
\centering
\includegraphics[width=\linewidth,keepaspectratio]{../../artifacts/figures/manuscript/artifact_benchmark/artifact_benchmark_eval_window41_Fp1_decomposition.png}
\caption{Per-channel decomposition for the Artifact Benchmark (validation window~41, channel~Fp1). Rows show raw versus MRANC-denoised overlay; EOG, EMG, ECG, and baseline noise stems; and normalized multi-scale attention weights.}\label{fig:ab-fp1-w41}
\end{figure*}

\begin{figure*}[!t]
\centering
\includegraphics[width=\linewidth,keepaspectratio]{../../artifacts/figures/manuscript/artifact_benchmark/artifact_benchmark_eval_window41_O1_decomposition.png}
\caption{Per-channel decomposition for the Artifact Benchmark (validation window~41, channel~O1). Rows show raw versus MRANC-denoised overlay; EOG, EMG, ECG, and baseline noise stems; and normalized multi-scale attention weights.}\label{fig:ab-o1-w41}
\end{figure*}

\begin{figure*}[!t]
\centering
\includegraphics[width=\linewidth,keepaspectratio]{../../artifacts/figures/manuscript/artifact_benchmark/artifact_benchmark_denoising_fidelity.png}
\caption{Denoising fidelity summary for the Artifact Benchmark.}\label{fig:ab-fidelity}
\end{figure*}

\end{document}
